\documentclass{../yalumba}

\title{Continuous Stability Weighting for\\Reference-Free Relatedness Detection:\\Module Clustering Stability v4}
\author{Ajax Davis}
\date{April 2026}

\setabstract{%
We present version~4 of Module Clustering Stability, a fundamental redesign of
the bootstrap resampling approach to reference-free relatedness detection.  Where
v3 treated stability as a binary filter---retaining modules appearing in $\geq 3$
of 5 bootstrap iterations and discarding the rest---v4 reconceives stability as
a \emph{continuous weight}.  Each motif receives a stability score
$\sigma \in [0, 1]$ representing the fraction of $B = 8$ bootstrap iterations
(at 70\% subsample) in which it participated in any module.  Pairwise relatedness
is scored by a two-component function: a generalized rarity-weighted Jaccard
index $\sum \min(\sigma_A, \sigma_B) \cdot r \,/\, \sum \max(\sigma_A, \sigma_B)
\cdot r$ and a rarity-weighted stability product
$\sum \sigma_A \sigma_B \cdot r \,/\, \sum \max(\sigma_A, \sigma_B)^2 \cdot r$,
where rarity $r = \log(N / \text{count})$ upweights motifs appearing in fewer
samples.  We systematically explore six weight configurations across three
second-component choices (supportCosine, rarityCosine, exclusiveRatio,
pure~genJaccard, prodRatio), achieving weakest gaps of $-3.98\%$ to $-1.04\%$.
The best configuration---genJaccard~$0.35$ + exclusiveRatio~$0.65$---reduces the
weakest gap from $-5.79\%$ (v3) to $-1.04\%$ but still fails threshold-based
classification.  Diagnosis reveals a structural confound: bootstrap stability
profiles correlate with sequencing coverage, not relatedness.  Samples with
similar depth produce similar stability fingerprints regardless of kinship,
a confound that cannot be resolved by weight tuning.  With Module
Persistence~v3 ($+3.36\%$ separation, $+0.13\%$ gap) and Coalition Transfer~v3
($+1.97\%$, $+0.06\%$ gap) both passing the CEPH~1463 pedigree test, Module
Stability remains the sole failing symbiogenesis algorithm, confirming that
the module identity problem---related individuals produce related but
non-identical modules---requires fuzzy matching or spectral methods rather than
improved weighting.}

\begin{document}
\maketitle

\tableofcontents
\newpage

% ════════════════════════════════════════════════════════════════════
\section{Introduction}
\label{sec:intro}

\subsection{From Binary to Continuous Stability}

The original Module Clustering Stability algorithm
(v3)~\cite{efron1979bootstrap, hennig2007cluster} treated bootstrap stability as
a binary predicate: a module either survived the threshold ($\sigma \geq 0.6$)
or was discarded entirely.  With $B = 5$ bootstrap iterations at 60\%
subsampling, stability scores could only take values in $\{0, 0.2, 0.4, 0.6,
0.8, 1.0\}$---six discrete levels, three of which (0--0.4) mapped to
``unstable'' and were zeroed out.

\ymarginnote{Binary stability wastes half the information.}

This binary approach had a clear limitation: it discarded information.  A motif
appearing in 2 of 5 iterations ($\sigma = 0.4$) was treated identically to one
appearing in 0 iterations ($\sigma = 0$), despite the former showing partial
reproducibility.  The scoring function then operated on a severely reduced
feature space: only 50--104 stable modules per sample out of thousands.

Version~4 addresses this by abandoning binary stability entirely.  Instead,
every motif that appears in \emph{any} bootstrap iteration receives a continuous
stability score $\sigma \in (0, 1]$, and this score is used directly as a
weight in the comparison function.  Motifs with high stability contribute more
to the pairwise score; motifs with low stability contribute less, but are not
silenced.  The threshold is replaced by a continuous decay: noise is suppressed
by multiplication rather than truncation.

\subsection{Increased Bootstrap Resolution}

Alongside the continuous weighting redesign, v4 increases the bootstrap count
from $B = 5$ to $B = 8$ and raises the subsample fraction from 60\% to 70\%.

\ymarginnote{$B = 8$ yields 9 distinct stability levels.}

With $B = 8$, stability scores take values in $\{0.125, 0.25, 0.375, 0.5,
0.625, 0.75, 0.875, 1.0\}$---eight non-zero levels, compared to five in v3.
Combined with the continuous weighting, this provides substantially finer
granularity for the scoring function.

The subsample fraction increase from 60\% to 70\% retains more data per
iteration, reducing the variance of module extraction while still providing
meaningful perturbation.  Each iteration removes 30\% of reads, sufficient to
disrupt noise-driven modules while preserving genuine structures.

\subsection{Context: The Symbiogenesis Program}

Module Clustering Stability~v4 is the latest iteration of the third algorithm
in the symbiogenesis family of reference-free relatedness detection
methods~\cite{margulis1998}.  The family comprises:

\begin{enumerate}
  \item \textbf{Module Persistence v3}: rarity-weighted Jaccard + support cosine
    on shared modules.  \emph{Separation: $+3.36\%$, weakest gap: $+0.13\%$.
    Status: PASSES.}
  \item \textbf{Coalition Transfer v3}: exclusive-ratio scoring of module
    coalitions (sets of modules co-occurring on reads).
    \emph{Separation: $+1.97\%$, weakest gap: $+0.06\%$. Status: PASSES.}
  \item \textbf{Module Clustering Stability v4}: continuous stability-weighted
    rarity scoring.  \emph{This paper.}
\end{enumerate}

Both Module Persistence~v3 and Coalition Transfer~v3 have achieved positive
weakest gaps on the CEPH~1463 pedigree, meaning they can separate all
parent--child pairs from all unrelated pairs.  Module Stability is the sole
remaining algorithm that fails this test.  The present paper documents the v4
redesign, the systematic weight search, and the structural diagnosis of why
weight tuning alone cannot resolve the failure.


% ════════════════════════════════════════════════════════════════════
\section{Methods}
\label{sec:methods}

\subsection{Dataset}

We evaluate on the CEPH~1463 three-generation pedigree~\cite{ceph1463,
zook2019giab, 1000genomes2015}, the standard benchmark throughout the yalumba
project.

\begin{pedigree}
Generation 1 (Grandparents):
  Paternal: NA12891 (Pat.GF) + NA12892 (Pat.GM)
  Maternal: NA12889 (Mat.GF) + NA12890 (Mat.GM)

Generation 2 (Parents):
  NA12877 (Father) = child of NA12891 + NA12892
  NA12878 (Mother) = child of NA12889 + NA12890
\end{pedigree}

This yields 10 pairwise comparisons: 4 parent--child, 2 spousal, and 4 in-law
or cross-lineage pairs.  Input is raw FASTQ at $\sim$30$\times$ coverage.
No alignment, variant calling, or phasing is performed.

\subsection{Bootstrap Resampling Protocol (v4)}

The v4 protocol differs from v3 in three parameters:

\begin{table}[H]
  \centering
  \tablestyle
  \caption{Bootstrap parameter comparison between v3 and v4.}
  \label{tab:params}
  \begin{tabular}{lcc}
    \toprule
    \rowcolor{table-header}
    \textcolor{white}{\textbf{Parameter}} &
    \textcolor{white}{\textbf{v3}} &
    \textcolor{white}{\textbf{v4}} \\
    \midrule
    Bootstrap iterations ($B$) & 5 & 8 \\
    Subsample fraction & 60\% & 70\% \\
    Stability usage & Binary ($\sigma \geq 0.6$) & Continuous weight \\
    Stability threshold ($\tau$) & 0.6 (3/5 minimum) & None (all $\sigma > 0$) \\
    Distinct $\sigma$ levels & 5 non-zero & 8 non-zero \\
    Max reads per sample & 30{,}000 & 30{,}000 \\
    \bottomrule
  \end{tabular}
\end{table}

For each sample, the protocol proceeds:

\begin{enumerate}
  \item Draw $\lfloor 0.7 \cdot N \rfloor$ reads uniformly without replacement.
  \item Execute the full module extraction pipeline on the subsample:
    \begin{itemize}
      \item Window size: 150\,bp
      \item K-mer size: $k = 15$
      \item Minimum support: 3 reads per k-mer
      \item Minimum cohesion: 0.25
    \end{itemize}
  \item Record which motifs (individual k-mers) participated in any discovered
    module, along with the maximum support observed for each motif.
  \item Repeat for $B = 8$ iterations total.
\end{enumerate}

\subsection{Motif-Level Stability Profiles}

\begin{definition}
The \textbf{stability score} of motif $m$ in sample $s$ is
\[
  \sigma(m, s) = \frac{|\{b : m \in \bigcup_{M \in \mathcal{M}_b(s)} M\}|}{B}
\]
where $\mathcal{M}_b(s)$ is the set of modules discovered in bootstrap iteration
$b$ for sample $s$, and $\bigcup M$ denotes the union of all motif members
across modules.
\end{definition}

\ymarginnote{v4 tracks motif stability, not module stability.}

A critical distinction from v3: the unit of stability tracking shifts from
\emph{modules} to \emph{motifs}.  In v3, a module was stable if its exact
k-mer set recurred across iterations.  In v4, each individual k-mer (motif)
receives its own stability score based on how many iterations it participated
in \emph{any} module.  This decouples stability from module identity, avoiding
the brittleness of exact set matching at the module level.

Each sample is represented by a \emph{motif stability profile}: a map from
motif hash to stability score.  The profile also accumulates cumulative module
support across iterations for ancillary analysis.

\subsection{Informative Filtering}

After profile construction, we apply population-level informative filtering.
A motif is considered ``present'' in a sample if its stability exceeds 0.25
(appears in at least 2 of 8 iterations).  Motifs present in \emph{all} samples
carry no discriminative power and are removed.

\begin{definition}
A motif is \textbf{informative} if the number of samples where
$\sigma(m, s) > 0.25$ is strictly less than $N$ (the total sample count).
\end{definition}

\subsection{Two-Component Scoring Function}

Given motif stability profiles for two samples $A$ and $B$, and the set
$\mathcal{I}$ of informative motifs, we compute two complementary scores.

\subsubsection{Component 1: Generalized Rarity-Weighted Jaccard}

The generalized Jaccard index~\cite{bray1957} extends the classical set Jaccard
to continuous-valued vectors:

\begin{equation}
  J_{\text{gen}}(A, B) =
  \frac{\displaystyle\sum_{m \in \mathcal{I}} \min(\sigma_A^m, \sigma_B^m) \cdot r_m}
       {\displaystyle\sum_{m \in \mathcal{I}} \max(\sigma_A^m, \sigma_B^m) \cdot r_m}
  \label{eq:gen-jaccard}
\end{equation}

where $\sigma_X^m = \sigma(m, X)$ is the stability score and the rarity weight
is:

\begin{equation}
  r_m = \log\!\left(\frac{N}{\text{count}(m)}\right)
  \label{eq:rarity}
\end{equation}

with $\text{count}(m)$ being the number of samples where $\sigma(m, \cdot) > 0.25$.
This follows TF-IDF principles~\cite{salton1988tfidf}: motifs appearing in fewer
samples receive higher rarity weight, amplifying the signal from lineage-specific
markers.

\ymarginnote{Rarity weighting amplifies lineage-specific signals.}

The generalized Jaccard naturally handles continuous stability: a motif with
$\sigma_A = 0.875$ and $\sigma_B = 0.125$ contributes
$\min(0.875, 0.125) / \max(0.875, 0.125) = 0.143$ before rarity weighting,
reflecting the low overlap in stability profiles.  In v3's binary scheme, one
motif would be ``stable'' and the other ``unstable,'' contributing zero.

\subsubsection{Component 2: Rarity-Weighted Stability Product}

The stability product rewards motifs that are simultaneously stable in both
samples:

\begin{equation}
  P(A, B) =
  \frac{\displaystyle\sum_{m \in \mathcal{I}} \sigma_A^m \cdot \sigma_B^m \cdot r_m}
       {\displaystyle\sum_{m \in \mathcal{I}} \max(\sigma_A^m, \sigma_B^m)^2 \cdot r_m}
  \label{eq:prod-ratio}
\end{equation}

This component captures a different signal than the generalized Jaccard.
Where $J_{\text{gen}}$ measures overlap of stability ranges, the product
$\sigma_A \cdot \sigma_B$ measures joint confidence.  A motif stable in both
samples ($0.875 \times 0.875 = 0.766$) receives far more weight than one
stable in only one ($0.875 \times 0.125 = 0.109$).  The normalization by
$\max^2$ ensures the ratio lies in $[0, 1]$.

\subsubsection{Composite Score}

The final score is a weighted combination:

\begin{equation}
  S(A, B) = w_1 \cdot J_{\text{gen}}(A, B) + w_2 \cdot P(A, B)
  \label{eq:composite}
\end{equation}

where $w_1 + w_2 = 1$.  We systematically explore six weight configurations
to find the optimal combination (Section~\ref{sec:weight-search}).


% ════════════════════════════════════════════════════════════════════
\section{Systematic Weight Search}
\label{sec:weight-search}

The v4 redesign introduces a new scoring architecture, and the optimal weights
are not known \emph{a priori}.  We systematically evaluate six configurations,
varying both the weight allocation and the choice of second component.

\subsection{Configurations Tested}

\begin{table}[H]
  \centering
  \tablestyle
  \caption{Weight configurations explored in the v4 iteration cycle.  Each row
    shows the two components, their weights, and the resulting weakest gap on
    the CEPH~1463 10-pair benchmark.  Negative weakest gap indicates
    classification failure.}
  \label{tab:weight-search}
  \begin{tabular}{lcccl}
    \toprule
    \rowcolor{table-header}
    \textcolor{white}{\textbf{Version}} &
    \textcolor{white}{\textbf{Component 1}} &
    \textcolor{white}{\textbf{Component 2}} &
    \textcolor{white}{\textbf{Weakest Gap}} &
    \textcolor{white}{\textbf{Notes}} \\
    \midrule
    v3 (baseline) & genJaccard (0.55) & supportCosine (0.45)
      & $-$2.34\% & Binary stability baseline \\
    v4a & genJaccard (0.55) & rarityCosine (0.45)
      & $-$3.24\% & Rarity-weighted cosine \\
    v4b & genJaccard (0.50) & exclusiveRatio (0.50)
      & $-$1.13\% & Exclusive motif fraction \\
    v4c & genJaccard (1.00) & --- (0.00)
      & $-$3.91\% & Pure generalized Jaccard \\
    v4d & genJaccard (0.35) & exclusiveRatio (0.65)
      & $\mathbf{-1.04\%}$ & \textbf{Best configuration} \\
    v4e & genJaccard (0.55) & prodRatio (0.45)
      & $-$3.98\% & Stability product ratio \\
    \bottomrule
  \end{tabular}
\end{table}

\subsection{Component Descriptions}

The six configurations explore three distinct second components alongside the
generalized Jaccard:

\begin{itemize}
  \item \textbf{supportCosine}: Cosine similarity of cumulative module support
    vectors across the informative motif space.  Captures whether two samples
    have similar support magnitudes, independent of stability.
  \item \textbf{rarityCosine}: Cosine similarity weighted by rarity.  Amplifies
    rare-motif contributions but proved to inflate in-law scores (pairs sharing
    coverage characteristics rather than ancestry).
  \item \textbf{exclusiveRatio}: Fraction of shared stable motifs that are
    exclusive to the pair (not found in other samples).  Rewards
    uniquely-shared genetic material.
  \item \textbf{prodRatio}: The stability product ratio defined in
    Equation~\ref{eq:prod-ratio}.  Rewards joint high-stability motifs.
\end{itemize}

\subsection{Weight Search Analysis}

\begin{keyfinding}
The exclusiveRatio component produces the best results.  Both configurations
using it (v4b: $-1.13\%$, v4d: $-1.04\%$) outperform all other second
components.  Increasing the exclusiveRatio weight from 0.50 to 0.65 improves
the weakest gap from $-1.13\%$ to $-1.04\%$, suggesting the exclusive-sharing
signal is the most informative available.
\end{keyfinding}

\begin{keyfinding}
Pure generalized Jaccard (v4c) performs poorly ($-3.91\%$ gap), demonstrating
that the Jaccard alone cannot separate related from unrelated pairs.  The
second component is essential, but no combination tested achieves a positive gap.
\end{keyfinding}

The prodRatio component (v4e: $-3.98\%$) performs worst, nearly matching
the pure Jaccard failure.  This indicates that joint stability---motifs stable
in both samples---is not a discriminative signal for relatedness on this dataset.
The reason becomes clear in Section~\ref{sec:diagnosis}: stability correlates
with coverage, not kinship.

\begin{limitation}
\textbf{All six configurations fail.}  The best weakest gap achieved ($-1.04\%$
for v4d) remains negative, meaning no threshold can cleanly separate related
from unrelated pairs.  The weight search space is exhausted within the current
scoring architecture.
\end{limitation}


% ════════════════════════════════════════════════════════════════════
\section{Results}
\label{sec:results}

We present detailed results for two configurations: the current implementation
(v4e: genJaccard~0.55 + prodRatio~0.45) and the best-performing configuration
(v4d: genJaccard~0.35 + exclusiveRatio~0.65).

\subsection{Per-Sample Stability Statistics (v4 Parameters)}

With the increased bootstrap count ($B = 8$) and subsample fraction (70\%),
the stability profiles differ from the v3 paper.

\begin{table}[H]
  \centering
  \tablestyle
  \caption{Per-sample motif stability statistics under v4 parameters ($B = 8$,
    70\% subsample).  ``Stable motifs'' counts those with $\sigma \geq 0.5$.
    All motifs with $\sigma > 0$ participate in scoring under continuous weighting.}
  \label{tab:v4-stability}
  \begin{tabular}{llrrrl}
    \toprule
    \rowcolor{table-header}
    \textcolor{white}{\textbf{Sample}} &
    \textcolor{white}{\textbf{Role}} &
    \textcolor{white}{\textbf{Total Motifs}} &
    \textcolor{white}{\textbf{$\sigma \geq 0.5$}} &
    \textcolor{white}{\textbf{Informative}} &
    \textcolor{white}{\textbf{Time (s)}} \\
    \midrule
    NA12889 & Mat.\ Grandfather & 8{,}793 & 50  & all & 77.6 \\
    NA12890 & Mat.\ Grandmother & 11{,}686 & 79  & all & 77.3 \\
    NA12891 & Pat.\ Grandfather & 12{,}068 & 91  & all & 65.3 \\
    NA12892 & Pat.\ Grandmother & 12{,}749 & 103 & all & 63.5 \\
    NA12877 & Father            & 12{,}999 & 85  & all & 60.6 \\
    NA12878 & Mother            & 11{,}842 & 104 & all & 60.6 \\
    \midrule
    \multicolumn{2}{l}{\textit{Total unique informative}} & --- & \textbf{503} & --- & --- \\
    \bottomrule
  \end{tabular}
\end{table}

\ymarginnote{Continuous weighting means all motifs contribute, not just the 503
at $\sigma \geq 0.5$.}

A critical difference from v3: under continuous weighting, the ``503 stable
modules'' are not the only features used.  Every motif with $\sigma > 0$ in at
least one sample contributes to the score, weighted by its stability.  The
high-stability motifs ($\sigma \geq 0.5$) contribute most, but the thousands
of lower-stability motifs collectively provide a secondary signal.

\subsection{Full Score Table (v4e: Current Implementation)}

\begin{table}[H]
  \centering
  \tablestyle
  \caption{Pairwise scores for v4e (genJaccard~0.55 + prodRatio~0.45), ranked
    by composite $S$.  Parent--child pairs marked with $\blacktriangleright$.}
  \label{tab:scores-v4e}
  \begin{tabular}{clccl}
    \toprule
    \rowcolor{table-header}
    \textcolor{white}{\textbf{Rank}} &
    \textcolor{white}{\textbf{Pair}} &
    \textcolor{white}{\textbf{$J_{\text{gen}}$}} &
    \textcolor{white}{\textbf{$S$}} &
    \textcolor{white}{\textbf{Type}} \\
    \midrule
    1 & $\blacktriangleright$ Mat.GM $\to$ Mother
      & --- & 0.2740 & Parent--child \\
    2 & Mat.GF $\leftrightarrow$ Father
      & --- & 0.2634 & In-law \\
    3 & $\blacktriangleright$ Mat.GF $\to$ Mother
      & --- & 0.2630 & Parent--child \\
    4 & $\blacktriangleright$ Pat.GM $\to$ Father
      & --- & 0.2625 & Parent--child \\
    5 & Spouses
      & --- & 0.2345 & Spousal \\
    6 & Pat.GM $\leftrightarrow$ Mat.GM
      & --- & 0.2315 & Unrelated \\
    7 & $\blacktriangleright$ Pat.GF $\to$ Father
      & --- & 0.2055 & Parent--child \\
    8 & Pat.GF $\leftrightarrow$ Mat.GF
      & --- & 0.1836 & Unrelated \\
    9 & Pat.GF $\leftrightarrow$ Mat.GM
      & --- & 0.1584 & Unrelated \\
    10 & Pat.GF $\leftrightarrow$ Mother
      & --- & 0.1418 & In-law \\
    \bottomrule
  \end{tabular}
\end{table}

\begin{itemize}
  \item \textbf{Mean separation}: $+4.90$ percentage points (parent--child
    avg.\ 0.2512 vs.\ unrelated avg.\ 0.2022).
  \item \textbf{Weakest gap}: $-3.98\%$.  The in-law pair
    Mat.GF~$\leftrightarrow$~Father (0.2634) exceeds the parent--child pair
    Pat.GF~$\to$~Father (0.2055).
\end{itemize}

\subsection{Best Configuration (v4d: genJaccard 0.35 + exclusiveRatio 0.65)}

The best-performing weight configuration (v4d) reduces the weakest gap to
$-1.04\%$, a $4.75$ percentage-point improvement over v4e.  The exclusive
ratio component, which measures the fraction of shared motifs unique to the
pair, provides the strongest discriminative signal among all tested components.

\begin{keyfinding}
\textbf{v4d achieves $-1.04\%$ weakest gap}---the closest any Module Stability
configuration has come to passing.  The exclusive-ratio signal is the most
informative second component, but 1.04 percentage points still separate the
method from the passing threshold.
\end{keyfinding}

\subsection{Comparison to v3}

\begin{table}[H]
  \centering
  \tablestyle
  \caption{Evolution from v3 to v4.  v3 results from the original paper;
    v4 results from the current implementation.}
  \label{tab:v3-v4}
  \begin{tabular}{lcccc}
    \toprule
    \rowcolor{table-header}
    \textcolor{white}{\textbf{Version}} &
    \textcolor{white}{\textbf{Scoring}} &
    \textcolor{white}{\textbf{Separation}} &
    \textcolor{white}{\textbf{Weakest Gap}} &
    \textcolor{white}{\textbf{Status}} \\
    \midrule
    v3 (paper) & 3-component, binary stability
      & $+$4.90\% & $-$5.79\% & Fails \\
    v4 (best, v4d) & 2-component, continuous stability
      & --- & $-$1.04\% & Fails \\
    v4 (current, v4e) & genJaccard + prodRatio
      & $+$4.90\% & $-$3.98\% & Fails \\
    \bottomrule
  \end{tabular}
\end{table}

The continuous weighting redesign achieves a substantial improvement in the
weakest gap ($-5.79\% \to -1.04\%$) for the best configuration, demonstrating
that the continuous approach extracts more signal from the stability profiles
than the binary filter.  However, the improvement is insufficient to cross
the zero threshold.

\subsection{Comparison to Sibling Algorithms}

\begin{table}[H]
  \centering
  \tablestyle
  \caption{All symbiogenesis algorithms at their latest versions, evaluated on
    the 6-member CEPH~1463 pedigree (10 pairs).}
  \label{tab:family}
  \begin{tabular}{lcccc}
    \toprule
    \rowcolor{table-header}
    \textcolor{white}{\textbf{Algorithm}} &
    \textcolor{white}{\textbf{Separation}} &
    \textcolor{white}{\textbf{Weakest Gap}} &
    \textcolor{white}{\textbf{Time (s)}} &
    \textcolor{white}{\textbf{Status}} \\
    \midrule
    Rare-Run P90 (baseline)
      & $+$583\% & $+$300\% & 72.5 & \textbf{Gold standard} \\
    Module Persistence v3
      & $+$3.36\% & $+$0.13\% & 1{,}085 & \textbf{PASSES} \\
    Coalition Transfer v3
      & $+$1.97\% & $+$0.06\% & 272 & \textbf{PASSES} \\
    \textbf{Module Stability v4 (best)}
      & --- & $\mathbf{-1.04\%}$ & 405 & Fails \\
    \bottomrule
  \end{tabular}
\end{table}

\begin{keyfinding}
Module Persistence~v3 and Coalition Transfer~v3 both achieve positive weakest
gaps, demonstrating that the CEPH~1463 benchmark is solvable within the
symbiogenesis framework.  Module Stability's continued failure ($-1.04\%$ at
best) indicates a problem specific to the bootstrap stability signal, not
the module paradigm itself.
\end{keyfinding}


% ════════════════════════════════════════════════════════════════════
\section{Structural Diagnosis}
\label{sec:diagnosis}

The systematic weight search (Section~\ref{sec:weight-search}) exhausted six
configurations across four second-component choices without achieving a positive
weakest gap.  This section diagnoses the structural cause.

\subsection{Bootstrap Stability Is Coverage-Correlated}

\begin{limitation}
\textbf{The core confound.}  Bootstrap stability profiles correlate with
sequencing coverage and quality, not with genealogical relatedness.  Samples
with similar depth produce similar stability fingerprints regardless of kinship.
This confound is structural---it cannot be resolved by weight tuning.
\end{limitation}

The mechanism is straightforward.  K-mer module extraction depends on read
coverage: a region with higher coverage produces more reads, which increases
the support for k-mers in that region, which makes the resulting modules more
robust to subsampling.  When we bootstrap by removing 30\% of reads, the
survival probability of a module is determined primarily by its original
support level, which is a function of coverage.

\ymarginnote{Coverage determines which modules survive perturbation.}

Two samples sequenced at similar depth will have similar coverage distributions
across the genome.  The same genomic regions will have high coverage in both
samples, producing modules with similar support levels and therefore similar
bootstrap stability scores.  The stability \emph{profile}---the vector of
per-motif stability scores---becomes a fingerprint of the coverage
distribution, not the genomic content.

\subsection{The NA12890/NA12892 Confound}

The confound is most visible in the pair NA12890 (Mat.\ Grandmother) and
NA12892 (Pat.\ Grandmother).  These two individuals are genealogically
unrelated---they married into opposite sides of the pedigree.  Yet they have
similar sequencing depth and quality in the CEPH~1463 dataset, producing
similar total module counts (11{,}686 vs.\ 12{,}749) and similar stable motif
counts (79 vs.\ 103).

\ymarginnote{Unrelated grandmothers with similar coverage produce similar
stability profiles.}

Under continuous stability weighting, this coverage similarity translates
directly into score similarity.  The generalized Jaccard and stability product
both reward motifs with similar $\sigma$ values across samples, and coverage
correlation ensures many motifs have similar $\sigma$ in NA12890 and NA12892.

The same confound affects the failing in-law pair
Mat.GF~$\leftrightarrow$~Father (NA12889 $\leftrightarrow$ NA12877).  Despite
being genealogically unrelated, these samples share enough coverage
characteristics to produce an inflated score (0.2634 in v4e) that exceeds the
true parent--child pair Pat.GF~$\to$~Father (0.2055).

\subsection{Why Weight Tuning Cannot Fix This}

The coverage confound operates \emph{within} every component of the scoring
function:

\begin{enumerate}
  \item \textbf{Generalized Jaccard} ($J_{\text{gen}}$): the min/max ratio
    of stability scores is inflated when two samples have similar coverage,
    because similar coverage $\Rightarrow$ similar $\sigma$ values
    $\Rightarrow$ min $\approx$ max $\Rightarrow$ high ratio.
  \item \textbf{Stability product} ($P$): the product $\sigma_A \cdot \sigma_B$
    is maximized when both values are high, which occurs for high-coverage
    motifs in high-coverage samples---independent of relatedness.
  \item \textbf{Rarity weighting}: rarity ($r_m = \log(N/\text{count})$)
    mitigates universal motifs but does not address within-rarity-class
    coverage correlation.
  \item \textbf{Exclusive ratio}: the best-performing component partially
    avoids the confound by requiring motifs to be unique to the pair, but
    coverage similarity between unrelated samples can produce spuriously
    ``exclusive'' motifs.
\end{enumerate}

No linear combination of these components can deconfound coverage from
relatedness, because the confound is embedded in the stability scores
themselves---the input to all components---not in the combination weights.

\subsection{The Pat.GF Vulnerability}

NA12891 (Pat.GF) remains the persistent weak link across all symbiogenesis
algorithms.  In Module Stability~v4, the Pat.GF~$\to$~Father pair scores
0.2055 (rank~7 of 10), the lowest of all parent--child pairs by a wide margin.

\begin{itemize}
  \item Pat.GF has 91 stable motifs ($\sigma \geq 0.5$) out of 12{,}068 total.
  \item Father has 85 stable motifs out of 12{,}999 total.
  \item The union of their high-stability motif sets contains 135 unique motifs
    with near-zero overlap.
  \item The count ratio ($12{,}068 / 12{,}999 = 0.928$) is reasonable, but the
    stability profile similarity is low because Pat.GF's coverage
    distribution differs from Father's.
\end{itemize}

\ymarginnote{Pat.GF confounds every symbiogenesis algorithm.}

The vulnerability is compounded by the in-law pair
Mat.GF~$\leftrightarrow$~Father, which benefits from a spuriously high score.
The combination---one parent--child pair scoring unusually low and one
unrelated pair scoring unusually high---creates the negative weakest gap.

\subsection{Module Identity Brittleness}

Beyond the coverage confound, a second structural issue persists from v3:
module identity is defined by exact k-mer set membership.  Even with continuous
stability weighting, the underlying motif hashes are exact---a motif either
matches or doesn't.  Related individuals who share a haplotype block will
produce k-mer modules from that block, but heterozygosity at even one SNP
within the block changes the k-mer composition, creating different motif
identities.

\begin{limitation}
\textbf{Exact motif identity cannot detect shared-but-heterozygous regions.}
A single SNP within a shared haplotype block changes the k-mer hash, making
the resulting motifs appear unrelated despite arising from inherited sequence.
\end{limitation}

This brittleness affects all module-based methods equally, but Module Stability
is most vulnerable because its scoring function operates entirely on motif-level
features.  Module Persistence and Coalition Transfer incorporate additional
structural signals (graph topology, coalition co-occurrence) that partially
compensate for identity brittleness.


% ════════════════════════════════════════════════════════════════════
\section{Discussion}
\label{sec:discussion}

\subsection{What Continuous Weighting Achieved}

The v3-to-v4 redesign demonstrates that continuous stability weighting is
strictly superior to binary filtering for this problem.  The best v4
configuration ($-1.04\%$) improves the weakest gap by 4.75 percentage points
over v3 ($-5.79\%$).  This improvement comes from two sources:

\begin{enumerate}
  \item \textbf{Information recovery.}  Binary filtering discards all motifs
    with $\sigma < 0.6$, losing the partial-stability signal.  Continuous
    weighting retains this signal, attenuated by the low $\sigma$ values.
    The aggregate contribution of thousands of low-stability motifs provides
    a secondary discriminant.
  \item \textbf{Finer granularity.}  With 8 non-zero stability levels instead
    of 5, the scoring function can distinguish between motifs at 62.5\% and
    87.5\% stability---a distinction that maps to biological relevance
    (moderately robust vs.\ highly robust structures).
\end{enumerate}

\ymarginnote{Continuous weighting is better, but not enough.}

However, continuous weighting does not address the fundamental coverage
confound.  The stability scores themselves are coverage-contaminated, and no
downstream weighting scheme can decontaminate them.

\subsection{The Limits of Weight Tuning}

The six-configuration weight search is, in retrospect, a study in diminishing
returns.  The first four configurations (v3, v4a, v4b, v4c) explored the
component space broadly, identifying exclusiveRatio as the best second
component.  The final two (v4d, v4e) refined the weights, gaining only 0.09
percentage points (v4b~$-1.13\%$ to v4d~$-1.04\%$).  Further tuning within
the existing architecture would yield marginal improvements at best.

The weight search demonstrates a general principle: \emph{when the features
are confounded, no classifier can deconfound them.}  The coverage signal is
embedded in the stability scores, and all components---Jaccard, cosine,
product, exclusive ratio---are computed from these scores.  Changing weights
redistutes the confound's influence but cannot eliminate it.

\subsection{Why Module Persistence and Coalition Transfer Pass}

The success of Module Persistence~v3 ($+0.13\%$ gap) and Coalition
Transfer~v3 ($+0.06\%$ gap) demonstrates that the coverage confound is not
universal to all module-based methods.  These algorithms succeed because they
incorporate signals beyond motif-level stability:

\begin{itemize}
  \item \textbf{Module Persistence} uses rarity-weighted Jaccard at the
    \emph{module} level (not motif level) combined with support cosine---a
    signal derived from module support patterns, not stability under
    perturbation.
  \item \textbf{Coalition Transfer} uses coalition structure: which modules
    co-occur on individual reads.  This captures physical linkage information
    (modules on the same read are from the same haplotype) that is orthogonal
    to coverage.
\end{itemize}

Module Stability's scoring function, by contrast, operates entirely on
per-motif stability scores---a single-dimensional feature space that is
saturated by the coverage confound.  The lesson is that dimensionality of
signal matters: methods that combine multiple orthogonal features can
average out confounds that dominate any single feature.

\subsection{Implications for Future Work}

The v4 diagnosis points to three concrete directions:

\begin{enumerate}
  \item \textbf{Coverage normalization.}  Explicitly model expected stability
    as a function of coverage, then score the \emph{residual} stability after
    subtracting the coverage effect.  This would require estimating per-region
    coverage from read depth, which is feasible but adds complexity.
  \item \textbf{Fuzzy motif matching.}  Replace exact motif identity with a
    similarity threshold: two motifs are ``equivalent'' if their k-mer Jaccard
    exceeds 0.8.  This would capture shared-but-heterozygous regions that
    exact matching misses.  The computational cost increases (pairwise motif
    comparison replaces hash lookup) but the biological validity improves
    substantially.
  \item \textbf{Hybrid stability + structure.}  Use bootstrap stability to
    identify high-confidence motifs, then score using Coalition Transfer-style
    co-occurrence structure.  This combines Stability's quality filtering with
    Coalition Transfer's orthogonal scoring signal.
\end{enumerate}

\subsection{Connection to Classical Bootstrap Theory}

The failure of Module Stability has an instructive connection to classical
bootstrap theory~\cite{hall1992bootstrap}.  The bootstrap is consistent for
statistics that are smooth functions of the empirical distribution.  Stability
scores satisfy this requirement: they are the frequency of a binary event
(module appearance) across bootstrap replicates, and the law of large numbers
ensures convergence to the true ``appearance probability'' as $B \to \infty$.

\ymarginnote{The bootstrap correctly estimates what it measures---but what it
measures is not what we need.}

The problem is not that the bootstrap estimate is wrong---it correctly
estimates the probability that a motif will appear in a module under random
subsampling.  The problem is that this probability is determined primarily by
coverage, not by the biological property (shared ancestry) that we want to
detect.  The bootstrap faithfully reports a confounded quantity.

This is a cautionary tale for the application of resampling methods in
genomics: bootstrap stability of a derived feature (k-mer modules) measures
the robustness of the \emph{derivation process}, not the biological relevance
of the derived feature.  When the derivation process is coverage-dependent,
stability inherits the coverage dependence.


% ════════════════════════════════════════════════════════════════════
\section{Computational Profile}
\label{sec:computational}

\subsection{Timing}

\begin{table}[H]
  \centering
  \tablestyle
  \caption{Timing breakdown for Module Stability~v4 ($B = 8$, 70\% subsample,
    6 samples).}
  \label{tab:timing}
  \begin{tabular}{lr}
    \toprule
    \rowcolor{table-header}
    \textcolor{white}{\textbf{Phase}} &
    \textcolor{white}{\textbf{Time}} \\
    \midrule
    Bootstrap extraction (6 samples $\times$ 8 iterations) & $\sim$405\,s \\
    \quad Per sample (average) & 67.5\,s \\
    \quad Per iteration (average) & 8.4\,s \\
    Pairwise comparison (10 pairs) & $<$ 0.01\,s \\
    \midrule
    \textbf{Total} & $\sim$\textbf{405\,s} \\
    \bottomrule
  \end{tabular}
\end{table}

\ymarginnote{8 iterations per sample but each processes 70\% of reads (vs.\
60\% in v3), yielding similar per-sample wall time.}

The 60\% increase in bootstrap iterations ($5 \to 8$) is offset by lower
per-iteration cost due to the 70\% subsample retaining more coherent module
structures that extract faster.  Total preparation time remains comparable
to v3 ($\sim$405\,s).

\subsection{Memory}

Heap delta remains near zero ($\sim$0\,MB net).  Each bootstrap iteration's
data structures (module lists, co-occurrence graphs) are garbage-collected
before the next iteration begins.  Only the motif stability profile
(a \texttt{Map<number, number>}) persists across iterations, requiring
negligible memory ($\sim$1\,MB per sample for $\sim$12{,}000 motif entries).

This memory efficiency contrasts sharply with Module Persistence~v3
(4{,}806\,MB peak heap), confirming that bootstrap approaches trade I/O
cost (multiple file reads) for memory efficiency (no persistent data
structures between iterations).


% ════════════════════════════════════════════════════════════════════
\section{Limitations}
\label{sec:limitations}

\begin{limitation}
\textbf{Coverage confound is structural.}  Bootstrap stability correlates
with sequencing coverage, contaminating all downstream scoring.  No weight
tuning can resolve this within the current architecture.
\end{limitation}

\begin{limitation}
\textbf{Exact motif identity.}  The hash-based motif comparison requires
bit-exact k-mer set membership.  Heterozygous SNPs within shared haplotype
blocks produce distinct motif hashes, preventing detection of inherited
similarity.
\end{limitation}

\begin{limitation}
\textbf{Weight search is grid-based.}  Only six configurations were tested.
A continuous optimization (gradient-free, given the non-differentiable
ranking metrics) might find better weights, but the structural confound
limits the achievable improvement.
\end{limitation}

\begin{limitation}
\textbf{Single dataset.}  All results are on CEPH~1463 (6 European-ancestry
individuals, 10 pairs).  Generalization to other ancestries, larger cohorts,
and heterogeneous sequencing depths is untested.
\end{limitation}

\begin{limitation}
\textbf{No coverage normalization.}  The method does not model or correct
for coverage-dependent stability.  Production use would require explicit
normalization.
\end{limitation}

\begin{limitation}
\textbf{Negative weakest gap persists.}  The best configuration ($-1.04\%$)
remains negative, precluding threshold-based classification.  Module Stability
cannot be used as a standalone relatedness detector in its current form.
\end{limitation}


% ════════════════════════════════════════════════════════════════════
\section{Conclusion}
\label{sec:conclusion}

Module Clustering Stability~v4 represents a principled redesign of the
bootstrap resampling approach to reference-free relatedness detection.  The
shift from binary stability filtering to continuous stability weighting, the
increase from $B = 5$ to $B = 8$ bootstrap iterations, and the systematic
exploration of six weight configurations collectively improve the weakest gap
from $-5.79\%$ (v3) to $-1.04\%$ (v4d)---a 4.75 percentage-point gain.

\ymarginnote{Better architecture, same fundamental limitation.}

Yet the improvement is insufficient.  The best configuration still fails
threshold-based classification, and the diagnosis is definitive: bootstrap
stability profiles are structurally correlated with sequencing coverage,
embedding a confound that no linear scoring function can resolve.

The v4 iteration closes the weight-tuning chapter for Module Stability.
Further progress requires architectural changes: coverage normalization to
deconfound stability from depth, fuzzy motif matching to detect
shared-but-heterozygous regions, or hybrid approaches that combine stability's
quality filtering with the orthogonal scoring signals used by Module
Persistence and Coalition Transfer.

\begin{pullquote}
The bootstrap correctly measures stability.  The problem is that stability
measures coverage, not kinship.
\end{pullquote}

With Module Persistence~v3 ($+0.13\%$ gap) and Coalition Transfer~v3
($+0.06\%$ gap) both passing the CEPH~1463 benchmark, the symbiogenesis
program has demonstrated that module-based reference-free relatedness detection
is feasible.  Module Stability's contribution is diagnostic: it identifies the
boundary between what weight tuning can achieve and what requires new signal
sources.  The Rare-Run~P90 baseline ($+583\%$ separation, $+300\%$ positive
gap, 72.5\,s) remains the unchallenged gold standard, but the symbiogenesis
family's incremental progress---from $-5.79\%$ to $+0.13\%$ across three
algorithms and twelve iterations---demonstrates that module-based approaches
are converging on viable relatedness detection.


% ════════════════════════════════════════════════════════════════════
% References
% ════════════════════════════════════════════════════════════════════

\bibliographystyle{plain}
\bibliography{../references}

\end{document}
